\section{Introduction}

Several AI safety proposals are premised on AI-driven formal verification (FV); this is a core hypothesis of the ARIA Safeguarding AI program and the related Guaranteed Safe AI agenda \cite{dalrymple2024guaranteed}. If FV is to play this role, it must scale to the complexities of real-world code, but today we have little evidence that such capabilities are possible. Relevant benchmarks are mostly small datasets crafted by hand~\cite{loughridge2024dafnybench}, larger datasets extracted from expert-driven verification projects~\cite{chakraborty2024neural}, or are focused on advanced mathematics rather than engineering~\cite{glazer2024frontiermath}. None of these are likely to be representative of the formal verification tasks required to verify real-world AI-generated software at scale.

We generate a new benchmark, \FVSpec, from public uses of property-based testing (PBT). While FV itself is rarely used in software development, PBT is a 'close sibling' technology which has seen significant adoption. In both FV and PBT, engineers write \emph{specifications}: logical properties that the software must obey. For example, we might require that a Python function \texttt{isort(lst)} always returns a sorted permutation of its input \texttt{lst}. The difference lies in how this property is checked. In FV, the property is proved using mathematical techniques---this results in almost-perfect confidence, but requires proof engineering by highly expert teams \cite{dodds2022formally}. In PBT, the code is randomly tested---e.g. for random values of x. This is a push-button process with no proof engineering.

PBT is typically much cheaper than FV, and as a result has seen wider adoption. 
In fact, PBTs are typically seen as the first step toward FV. 
For example, in the interactive theorem prover ACL2s, any unproven theorem is automatically treated as a PBT (and tested against randomly generated inputs) until the author manages to supply a proof. Thus the idea of translating PBTs to theorems for FV is natural and, to an extent, industry-standard. We build on this intuition. First, we scrape $\approx$20k permissively-licensed PBTs written in Python's Hypothesis framework~\cite{maciver2019hypothesis}. Each PBT in our dataset can be seen as a small theorem statement (or \emph{ansatz}) about a piece of Python code, which the author would like to prove or, failing that, at least test. To make it possible for an AI to prove or refute these theorems, we then translate both code and PBTs to the theorem prover Lean, which is becoming the de-facto standard in AI formal verification. This approach was pioneered by Dougherty and Mehta on the FVAPPS verification benchmark \cite{dougherty2025proving}. In FVAPPS, source programs were taken from the APPS benchmark by Hendrycks et al \cite{hendrycks2021apps}, translated to PBTs, and then lifted to theorems. We apply the same process at scale to PBTs found in the wild.

We use AI-driven translation to generate Lean programs and theorems at scale. LLMs are quite good at program translation, and FVAPPS has shown that this approach is practical for constructing Lean theorems. 
Although in some cases the translations may be imperfect, lessening the importance of any proof or refutation of the resultant theorem from the perspective of the author of the originating PBT, it is nevertheless the case that we produce high-quality and realistic FV challenges for AI benchmarking.
The result is tens of thousands of verification challenges, each consisting of (1) a program and (2) a desired specification, both written in Lean. Unlike existing FV benchmarks, each problem corresponds to properties of real-world software written by engineers with no guaranteed formal verification experience.

\TODO: Summary of results.